% ---------------------------------------------------------------------------
% 1 Introduction
% Source: context/01_scenario_and_motivation.md, context/02_research_timeline.md
%
% This paper asks a feasibility question and answers it. The shape of the section is therefore:
% tension -> the question -> why the question is hard -> how we answer it -> the two verdicts.
% ---------------------------------------------------------------------------

\section{Introduction}

Human genomic sequence carries an unusually persistent confidentiality risk. It can identify its
contributor and, through familial relationships, reveal information about people who never supplied
a sample; re-identification has been demonstrated even when conventional identifiers were
removed~\cite{gymrek2013,erlich2018}. A disclosed genome also cannot be replaced in the manner of a
password or access token. European data-protection law accordingly treats genetic data as a special
category~\cite{gdpr}. The same sequence is the input on which genomic foundation models derive
their value: DNAGPT serves classification, regression, and generation tasks through a common
transformer backbone~\cite{dnagpt2023}. When that model is served rather than distributed, a data
owner without accelerator hardware must otherwise disclose the sequence to the provider or forgo
the inference. Nor is handing over an embedding a way out: per-token embeddings from DNA foundation
models have been shown to permit near-perfect reconstruction of the underlying
sequence~\cite{dnaembeddinginversion2026}. Homomorphic encryption offers a route for the provider to
compute without receiving the genomic values in plaintext.

DNAGPT's public release makes the experiment auditable and also makes this particular artifact
available for local execution. We use it as a reproducible surrogate for a genomic model offered
only as a service by policy or contract. Under that premise, client-assisted evaluation gives a
CPU-equipped data owner access to the served model without disclosing the genome.

Against this deployment premise, we ask: can inference with a released genomic foundation model be
executed under homomorphic encryption, and where would correctness, performance, or memory prevent
a complete encrypted deployment? The question contains two verdicts that should not be collapsed.
Feasibility asks whether the released computation can be executed faithfully within the available
cryptographic and memory budget. Practicality asks whether its resource cost supports a target use.
A correct but slow encrypted path is therefore evidence: it resolves the first question while
locating the work needed to change the second.

The technical difficulty is concentrated where a transformer departs from affine arithmetic. CKKS
supports approximate packed arithmetic and maps naturally to dense projections, rotations, and
accumulations~\cite{ckks2017}; LayerNorm, softmax, and GELU are not native CKKS operations. A
non-interactive evaluator must approximate them with polynomial circuits, lengthening the
multiplicative path and increasing the context, bootstrapping, and evaluation-key material that
must be resident. Moving the exact functions to the key holder shortens those paths and resets the
depth budget, but introduces interaction and a different threat model. Earlier systems establish
both directions: THOR demonstrates non-interactive homomorphic transformer inference; THE-X and
Safhire use interactive workflows in which the client evaluates nonlinear
functions~\cite{thor,thex2022,safhire2025}. The division of labor is therefore established prior
art. What has not been reported is an auditable account of the numerical behavior, memory,
interaction, and executed work that results when it is applied end to end to a released genomic
transformer, at that model's own task prompt, with its weights and nonlinear formulas unmodified.

That choice buys exactness and depth. Evaluating LayerNorm, softmax, and GELU at the key holder
keeps the released formulas intact, so no polynomial approximation of a nonlinearity enters the
model, and it bounds multiplicative depth by the longest encrypted segment rather than by network
depth. The price is an online key holder and a protocol that is not the fastest available:
non-interactive systems report lower cost for transformer inference, and Section~\ref{sec:related}
treats those figures as cross-paper context rather than as a controlled benchmark.

We begin by establishing that the target computation is meaningful and independently reproducible.
The released DNAGPT paths are evaluated across genomic-signal recognition, promoter and splice-site
classification, and mRNA abundance regression, with task outputs frozen before encrypted work.
The validation records the released float32 path, checks an independent NumPy float64 lineage
against a float64-cast PyTorch shadow at $2\times10^{-5}$ for every released block and the task
head, and then compares encrypted execution with the NumPy reference. We build two CKKS designs: a
non-interactive circuit with polynomial nonlinearities, and a client-assisted circuit whose fixed,
non-adaptive boundaries evaluate the released nonlinearities exactly. Instrumented evaluations
record where each design stops in correctness and memory, while runtime assertions preserve the
complete cryptographic operation schedule throughout the systems campaign.

The results separate three candidate barriers and resolve two of them. Correctness holds: the
client-assisted design evaluates all twelve released blocks and the task head at the $103$-token
prompt, and for the one input evaluated it reproduces the frozen plaintext label with
$8.56\times10^{-9}$ relative margin error against a $4\times10^{-2}$ tolerance fixed before
evaluation. Accelerator memory holds as well, and for a structural reason: because multiplicative
depth is bounded by the longest segment between client boundaries instead of by network depth, it
stays at $13$ for the complete model, device-wide GPU memory used reaches $9839$~MiB without growth
across the twelve blocks, and no homomorphic bootstrap is required at any point. Bounded caching
separately prevents encoded plaintexts from accumulating across host-side stages. The evaluated
non-interactive design, by contrast, completes one released-weight block but exceeds the tested
$80$~GB envelope while loading its composition context and evaluation keys, placing that barrier in
the nonlinear schedule rather than in the model's dense maps. Cost is the barrier that remains: one
encrypted classification takes $1.86$~hours on an A100 GPU node, which forecloses interactive use
in the evaluated form, and sampled GPU utilization peaks at $51\%$, locating headroom without
attributing it. These memory and timing figures are specific to the evaluated backend, packing, and
parameters.

The paper makes the following contributions:

\begin{itemize}
  \item \textbf{Complete-model feasibility.} All twelve released DNAGPT blocks and the task head
        execute at the $103$-token task prompt from encrypted embedded vectors, with unmodified
        released weights and nonlinear formulas and without homomorphic bootstrap; the one
        evaluated prompt reproduces the reference label. To our knowledge this is the first
        complete forward pass of a released, pretrained genomic foundation model evaluated under
        CKKS; prior homomorphic transformer evaluations over biological sequence data have targeted
        small, purpose-trained classifiers~\cite{idash24descriptor,encrypteddashformer}.
  \item \textbf{Model-faithful, reproducible evaluation.} The protocol, packing, checked operation
        schedule, predeclared tolerance, and three-link float32--float64--encrypted validation
        chain make the complete execution independently inspectable.
  \item \textbf{Measured systems anatomy and scoped negative results.} We report complete-model
        cost, interaction, depth, and memory; the stage-bounded cache that makes host-side reuse
        executable; and the memory boundary reached by the evaluated non-interactive configuration.
\end{itemize}
